% !TeX program = xelatex
% UTF-8 transcription; numerical values and abbreviations follow the notes.
% pic01.jpg = 20261012_211920204.jpg (page 1)
% pic02.jpg = 20261012_211926338.jpg (page 2)
% pic03.jpg = 20261012_211931934.jpg (page 3)
% pic04.jpg = 20261012_211936588.jpg (page 4)
% Compile directly with XeLaTeX or input into a XePersian book.
\newif\ifParaclinicStandalone
\ifdefined\settextfont
  \ParaclinicStandalonefalse
\else
  \ParaclinicStandalonetrue
\fi
\ifParaclinicStandalone
  \documentclass[a4paper,12pt]{extbook}
  \usepackage[margin=2.5cm]{geometry}
  \usepackage{xcolor}
  \usepackage{amsmath}
  \usepackage{enumitem}
  \usepackage{xepersian}
  \settextfont{Amiri}
  \setlatintextfont{FreeSerif}
  \pagestyle{plain}
\fi
\ifParaclinicStandalone
  \begin{document}
\fi
\begingroup
\emergencystretch=3em
\newcommand{\PCen}[1]{\lr{#1}}
\newcommand{\PCmath}[1]{\lr{$#1$}}
\ifdefined\colorbox
  \newcommand{\PChl}[1]{\colorbox{yellow!35}{\strut #1}}
\else
  \newcommand{\PChl}[1]{\textbf{#1}}
\fi
\newcommand{\PCunclear}[1]{#1 [خوانش نامطمئن]}
\newcommand{\PCpage}[1]{%
  \clearpage\section*{صفحهٔ #1 ـ ۲ ـ پاراکلینیک در نورولوژی}%
  {\footnotesize\noindent\PCen{Date: / /}\quad
  \PCen{Sat. Sun. Mon. Tue. Thu. Wed. Fri.}\quad\PCen{Subject:}\par}\medskip}

\chapter{پاراکلینیک در نورولوژی}

\PCpage{۱}
\section*{کنتراندیکاسیون‌های \PCen{LP}}
\begin{itemize}
\item \PChl{عفونت موضعی} محل پانکشن: \PChl{تنها کنتراندیکاسیون مطلق}.
  \begin{itemize}
  \item پانکشن از یک فضا بالاتر یا پایین‌تر؛ \PCunclear{نکردن، بهترین}.
  \end{itemize}
\item شک به تودهٔ داخل جمجمه: ریسک هرنی ترانس‌تنتوریال یا تونسیلار مخچه‌ای.
  \begin{itemize}
  \item \PChl{سردرد یا ادم پاپی}: \PCen{RF} برای وجود ضایعهٔ فضاگیر.
  \item قبل از \PCen{LP}، \PChl{\PCen{CT} یا \PCen{MRI}}.
  \end{itemize}
\item شک به تودهٔ نخاع: باید قبلش تصویربرداری کرد.
\item اختلال انعقادی: \PChl{\PCmath{plt < 20{,}000\text{--}50{,}000}}.
  \begin{itemize}
  \item اختلال عملکرد پلاکت (\PChl{الکلیسم، اورمی}).
  \item کمبود فاکتورها (\PChl{\PCmath{INR > 1.4}}، \PChl{\PCmath{PT > 14}}).
  \end{itemize}
\end{itemize}

\section*{عوارض \PCen{LP}}
\begin{itemize}
\item \PChl{سردرد: شایع‌ترین عارضه} / در \PCmath{\tfrac{1}{3}} / با ایستادن تشدید و با دراز کشیدن بهتر می‌شود.
  \begin{itemize}
  \item درمان با \PChl{\PCen{NSAID} یا کافئین}.
  \item در صورت مقاومت به درمان: یک \PCen{Patch} از خون بیمار در فضای اپیدورال محل نشت \PCen{CSF}.
  \end{itemize}
\item \PChl{تجویز مایعات زیاد} یا خواباندن بیمار در بستر به مدت ۱–۲ ساعت احتمال سردرد را \PChl{کاهش نمی‌دهد}.
\item خونریزی: جهت درمان اصلاح کوآگولوپاتی / گاهی خارج کردن لخته با جراحی.
  \begin{itemize}
  \item بررسی بیمار از نظر فشار بر نخاع، توسط هماتوم اپیدورال یا ساب‌دورال.
  \end{itemize}
\end{itemize}
\section*{\PCen{CSF} خونی}
\PCen{LP} تروماتیک / خونریزی ساب‌آراکنوئید (\PChl{\PCen{SAH}}).

\PCpage{۲}
\section*{موارد به نفع \PCen{LP} تروماتیک}
\begin{enumerate}[label=\protect\PCen{\alph*)}]
\item \PChl{فشار \PCen{CSF} طبیعی} (\PChl{\PCmath{100\text{--}180\ \mathrm{mmH_2O}}}).
\item به‌تدریج خون کم شود.
\item در نمونه‌های متوالی \PCen{RBC} کم شود.
\item اگر هماتوکریت بالا باشد، ۱–۲ \PCen{WBC} به ازای هر ۱۰۰۰ \PCen{RBC}.
\item اگر نمونهٔ \PCen{CSF} را بلافاصله سانتریفیوژ کنیم، موارد زیر به نفع:
  \begin{itemize}
  \item اگر بخش فوقانی نمونه بی‌رنگ باشد: \PCen{LP} تروماتیک.
  \item وجود \PCunclear{اریتروکرومی} در ابتدا و سپس زانتوکرومی پس از ۲۴ ساعت: \PCen{SAH}.
  \item سایر علل زانتوکرومی \PCen{CSF}:
    \begin{itemize}
    \item \PCunclear{\PCmath{Bil > 4\text{--}9\ \mathrm{mg/dl}}}.
    \item \PCmath{pro > 150\ \mathrm{mg/dl}}.
    \item وجود پیگمان کاروتن.
    \end{itemize}
  \end{itemize}
% The handwritten list jumps from e to g; its lettering is preserved.
\setcounter{enumi}{6}
\item لخته شدن مایع \PCen{CSF}: به نفع \PCen{LP} تروماتیک است.
\end{enumerate}

\section*{آنالیز \PCen{CSF}}
\begin{itemize}
\item فشار \PCen{CSF}: در بالغین \PChl{\PCmath{100\text{--}180\ \mathrm{mmH_2O}}} / در کودکان \PChl{\PCmath{30\text{--}60\ \mathrm{mmH_2O}}}.
\item \PChl{فشار \PCmath{>200\ \mathrm{mmH_2O}}}: \PCmath{ICP\uparrow}.
  \begin{itemize}
  \item توده‌های اینتراکرانیال، مننگوانسفالیت، \PCen{SAH}، سودوتومور سربری.
  \end{itemize}
\item ظاهر \PCen{CSF}: شفاف و بی‌رنگ.
\item اگر \PChl{\PCmath{WBC > 300/\mathrm{mm}^3}}: \PChl{کدر} می‌شود.
\item حداکثر \PChl{۵ لنفوسیت یا مونوسیت} در هر \PCmath{\mathrm{mm}^3}: \PChl{\PCen{PMN} نباید باشد}.
\item علل \PCmath{WBC\uparrow}: عفونت، واسکولیت، نئوپلاسم، التهاب ایمونولوژیک، واکنش به خون، مواد شیمیایی.
\end{itemize}

\PCpage{۳}
\section*{ادامهٔ آنالیز \PCen{CSF}}
\begin{itemize}
\item \PCen{pro} در \PCen{CSF}: در بالغین \PChl{\PCmath{25\text{--}45\ \mathrm{mg/dl}}} / در کودکان \PChl{\PCmath{<20\ \mathrm{mg/dl}}}.
\item علل \PCmath{pro\uparrow}: \PChl{مننژیت باکتریال} (\PChl{\PCmath{>500}})، \PChl{عفونت ویروسی \PCen{CSF}} (\PChl{\PCmath{\leq200}})، تومورهای اطراف طناب، گیلن‌باره، انسداد در جریان \PCen{CSF}.
\item علل \PCmath{pro\downarrow}: \PChl{هایپرتیروئیدی}، \PChl{\PCmath{ICP\downarrow}}.
\item \PCen{pro} خاص در \PCen{CSF}:
  \begin{itemize}
  \item تائو پروتئین: آلزایمر.
  \item \PChl{گاماگلوبولین: \PCen{MS}}.
  \item \PCen{14-3-3}: بیماری پریون.
  \end{itemize}
\item گلوکز در \PCen{CSF}: \PChl{\PCmath{\tfrac{2}{3}} غلظت گلوکز سرم} / ۴۰–۸۰ / \PChl{مقادیر \PCmath{<35} غیرطبیعی}.
  \begin{itemize}
  \item نمونهٔ \PCen{CSF} و خون باید همزمان و ناشتا گرفته شود و یا نمونهٔ سرم ۲–۴ ساعت قبل از نمونهٔ \PCen{CSF} گرفته شود.
  \end{itemize}
\item علل \PChl{گلوکز \PCmath{\downarrow} + پلئوسیتوز}:
  \begin{itemize}
  \item مننژیت چرکی، سلی یا قارچی.
  \item \PChl{مننژیت کارسینوماتوز}.
  \item \PChl{سارکوئیدوز}.
  \item \PChl{\PCen{SAH} در هفتهٔ اول}.
  \end{itemize}
\item عفونت‌های ویروسی سبب \PCmath{\downarrow} گلوکز نمی‌شوند.
\end{itemize}

\section*{\PCen{EEG}}
\begin{itemize}
\item مفیدترین تست تشخیصی در بیماران مشکوک به صرع.
\item کاربرد تشخیصی: وجود فعالیت اپی‌لپتی‌فرم: قویاً به نفع صرع.
  \begin{itemize}
  \item در ۲\% افراد طبیعی نیز وجود دارد.
  \end{itemize}
\item کاربرد در تقسیم‌بندی صرع، انتخاب درمان مناسب، تعیین پیش‌آگهی، تصمیم برای قطع دارو، تشخیص استاتوس اپیلپتیکوس بدون حرکات تشنجی.
\item \PChl{امواج فوکال آهسته} در \PCen{EEG} به نفع وجود \PChl{ضایعهٔ فوکال} است.
\end{itemize}

\PCpage{۴}
\section*{ادامهٔ \PCen{EEG}}
\begin{itemize}
\item \PCen{EEG} در اختلال سطح هوشیاری: تغییرات \PCen{EEG} با توجه به اتیولوژی، غیراختصاصی است ولی با شدت اختلال سطح هوشیاری مرتبط است.
  \begin{itemize}
  \item کاربرد: بررسی علل متابولیک، لوکالیزاسیون فوکال یا منتشر بودن ضایعه؛ بررسی وجود استاتوس بدون حرکات، ارزیابی پیش‌آگهی.
  \end{itemize}
\item \PCen{EEG} در بیمار مرگ مغزی ایزوالکتریک است.
\item امواج پریودیک در \PCen{EEG}:
  \begin{itemize}
  \item \PChl{پان‌انسفالیت اسکلروزان تحت حاد}، \PCen{SASPE}.
  \item \PChl{بیماری کروتزفلد ـ یاکوب}.
  \end{itemize}
\end{itemize}

\section*{\PCen{CT}}
\begin{itemize}
\item وضوح خون از ابتدای خونریزی در \PCen{CT} نسبت به \PCen{MRI} بهتر است.
\item کلسیم، چربی، استخوان مخصوصاً در کف جمجمه و پرده‌ها در \PCen{CT} بهتر دیده می‌شوند.
\item \PChl{در بیمار مشکوک به سکتهٔ مغزی حاد}، اولین اقدام تشخیصی \PChl{\PCen{CT} بدون کنتراست}.
  \begin{itemize}
  \item در موارد ایسکمیک تا \PCen{48h} اولیه چیزی دیده نمی‌شود.
  \end{itemize}
\end{itemize}

\section*{موارد برتری \PCen{MRI}}
\begin{itemize}
\item افتراق بهتر مادهٔ سفید و خاکستری.
\item ضایعات \PCunclear{کوچک}، ورتکس، حفرهٔ خلفی، نخاع، هیپوفیز.
% One word before "ورتکس" has been overwritten in the source.
\item تعیین پلاک‌های \PCen{MS}.
\item تعیین ضایعات مولد صرع (\PChl{دیسپلازی کورتکس}).
\item تشخیص سریع انفارکتوس مغز طی چند ساعت پس از آغاز و \PCunclear{مدت آن} (\PCen{DWI}).
\item تعیین سیگنال‌های غیرطبیعی مادهٔ سفید.
\item پاتولوژی‌های پارانشیم مغز و نخاع، تعیین اندازهٔ تومورها.
\item امکان تصویر برداشتن چندسطحی.
\item کشف ادم، التهاب \PCunclear{مفصل} و آبسه.
\end{itemize}

\endgroup
\ifParaclinicStandalone
  \expandafter\enddocument
\fi
